\documentclass[11pt,a4paper]{article}
\usepackage{review-style}
% Standalone illustrative sample: not course content or an assessment requirement.
% Substantive proofs get Core idea + conditions + reconstructible named actions.
% Elementary support stays in 1--3 lines without compulsory scaffolding.
% No compulsory numerical example, lookup sheet, or full-course frontmatter.
\renewcommand{\reviewtitle}{Proof layout sample}
\begin{document}
% Elementary support: no Core idea, numbered Steps, or appended numeric check.
\reviewsection{Indicator expectation: elementary support}{indicator-support}
For an event $A$ on a probability space,
$\mathbb E[\mathbf 1_A]=1\cdot\mathbb P(A)+0\cdot\mathbb P(A^c)=\mathbb P(A)$.
% Assessment labels must come from the source; complexity controls layout.

\reviewsection{Minimum-norm solution: orthogonal decomposition}{minimum-norm}
\provenance{Constructed layout demonstration. Assessment status must come from the actual course.}
\begin{reviewmethod}{Minimum-norm solution}
\assumptions{$J\in\mathbb R^{m\times p}$, $q\in\mathbb R^m$, and $K=JJ^\top\succ0$.}
The unique minimum-norm solution of $Ju=q$ is $u_*=J^\top K^{-1}q$.
\end{reviewmethod}
\coreidea{Split each feasible solution into the candidate in the row space and an orthogonal null-space residual. Only the residual adds squared norm.}
\begin{enumerate}
\item \textbf{Establish feasibility.} Substitute the candidate using invertibility:
\[
Ju_*=JJ^\top K^{-1}q=KK^{-1}q=q.
\]
\item \textbf{Identify the residual.} For any feasible $u$, let $z=u-u_*$. Both solutions satisfy the constraint, so
\[
\begin{aligned}
Jz&=Ju-Ju_*=q-q=0,\\
u_*^\top z&=q^\top K^{-1}Jz=0.
\end{aligned}
\]
\item \textbf{Compare norms.} Orthogonality gives
\[
\|u\|^2=\|u_*\|^2+\|z\|^2\ge\|u_*\|^2.
\]
Equality holds exactly when $z=0$, proving uniqueness.
\end{enumerate}
\examtrap{If $K$ is singular, the displayed inverse formula is unavailable. A pseudoinverse treatment needs its own feasibility conditions. The original positive-definiteness assumption stays with the statement.}
% No arbitrary Example is appended to fill a per-result quota.
% Later uses point to the canonical argument instead of reprinting its proof.
The row-space decomposition in \hyperref[minimum-norm]{Minimum-norm solution}
(p.~\pageref{minimum-norm}) also explains why a nonzero feasible null-space
residual strictly increases the norm.
% Source-optional extensions may go at lecture end or in an authorized volume.
% Remove duplicate later-covered previews; do not move assessed APPLY out.
% This template is not evidence of contraction: measure the actual task output.
\end{document}
